\chapter{2016:Jean-Pierre Sauvage and Sir J. Fraser Stoddart and Bernard L. Feringa}

\label{chap:chemistry2016}

The Nobel Prize in Chemistry for the year 2016 was awarded jointly to Jean-Pierre Sauvage and Sir J. Fraser Stoddart and Bernard L. Feringa.

\textbf{Motivation:}

for the design and synthesis of molecular machines

\section{Jean-Pierre Sauvage}

\subsection*{Early Life and Education}

Jean-Pierre Sauvage was born in 1944 in Paris, France.

\subsection*{Career and Major Research}

Sauvage received his doctorate from the University of Strasbourg and became a director of research of the CNRS at the University of Strasbourg. In 1983 he achieved the first synthesis of a catenane, two interlocked molecular rings, by threading one ring through a template and joining the second ring around it.

\subsection*{Nobel Prize-winning Work}

The work for which Jean-Pierre Sauvage was awarded the Nobel Prize in Chemistry in 2016 was described by the Nobel Committee as: "for the design and synthesis of molecular machines".

Sauvage's template method made it possible to construct molecules whose parts are mechanically linked yet free to move relative to one another.

\subsection*{Significance and Impact}

Interlocked molecules such as catenanes are the basis of molecular machines, since their linked parts can be made to move in a controlled way.

\subsection*{Later Career and Honors}

Sauvage is professor emeritus at the University of Strasbourg, where he continues research on molecular knots and interlocked structures.

\subsection*{Legacy}

His synthesis of catenanes opened the field of mechanically interlocked molecules, now central to molecular nanotechnology.

\subsection*{Modern Developments}

Sauvage's synthesis of catenanes and rotaxanes, molecular structures in which rings are mechanically interlocked, founded the field of mechanically interlocked molecules and molecular machines. His work laid the groundwork for molecular motors, switches, and pumps that are now used in molecular nanotechnology, drug delivery, and the design of artificial molecular machines.

\subsection*{Selected Publications}

"Une Nouvelle Famille de Molécules: les Métallo-Caténanes" (Tetrahedron Letters, 1983)

\clearpage

\section{Sir J. Fraser Stoddart}

\subsection*{Early Life and Education}

Sir J. Fraser Stoddart was born in 1942 in Edinburgh, Scotland.

\subsection*{Career and Major Research}

Stoddart earned bachelor's and doctoral degrees from the University of Edinburgh, receiving his doctorate in 1966. He was a professor at Northwestern University and later at the University of California, Los Angeles, and the University of Hong Kong. He developed methods for making rotaxanes, rings threaded on molecular axles with stoppers at each end, and used them to build molecular switches.

\subsection*{Nobel Prize-winning Work}

The work for which Sir J. Fraser Stoddart was awarded the Nobel Prize in Chemistry in 2016 was described by the Nobel Committee as: "for the design and synthesis of molecular machines".

Stoddart showed that a ring in a rotaxane can be moved back and forth along its axle in response to a stimulus, producing a controllable molecular shuttle.

\subsection*{Significance and Impact}

His molecular shuttles and switches demonstrated that molecules could perform mechanical tasks on command, laying the groundwork for nanoscale machines.

\subsection*{Later Career and Honors}

Stoddart was chair professor of chemistry at the University of Hong Kong at the end of his career. He died on 2024-12-30 in Melbourne, Australia.

\subsection*{Legacy}

Stoddart's interlocked molecules underpin molecular electronics and have inspired proposals for molecular computers and drug-delivery systems.

\subsection*{Modern Developments}

Stoddart's synthesis of rotaxanes and his construction of molecular switches and motors, including a molecular elevator and a molecular chain, demonstrated that molecules can perform mechanical functions. His work founded the field of artificial molecular machines, now being developed for use in nanoelectronics, drug delivery, and adaptive materials.

\subsection*{Selected Publications}

"Molecular Meccano. 1. [2]Rotaxanes and a [2]Catenane Made to Order" (Journal of the American Chemical Society, 1992)

\clearpage

\section{Bernard L. Feringa}

\subsection*{Early Life and Education}

Bernard L. Feringa was born in 1951 in Barger-Compascuum, the Netherlands.

\subsection*{Career and Major Research}

Feringa received his doctorate from the University of Groningen in 1978 and became a professor there. In 1999 he built the first molecular motor driven by light, a molecule that rotates in one direction only when illuminated.

\subsection*{Nobel Prize-winning Work}

The work for which Bernard L. Feringa was awarded the Nobel Prize in Chemistry in 2016 was described by the Nobel Committee as: "for the design and synthesis of molecular machines".

Feringa's light-driven motor converts light energy into a continuous unidirectional rotation of one part of the molecule relative to the other.

\subsection*{Significance and Impact}

His molecular motors were the first molecules that could be made to rotate continuously in a controlled direction, a requirement for useful machines at the molecular scale.

\subsection*{Later Career and Honors}

Feringa continues his research on molecular motors and rotary machines at the University of Groningen.

\subsection*{Legacy}

Feringa's motors are being developed for applications such as responsive materials, and they established that molecular machines can do real work.

\subsection*{Modern Developments}

Feringa's design of the first molecular motor, a molecule that rotates continuously when illuminated, demonstrated that molecules can act as tiny engines. His motors are now used to power nanoscale machines, create responsive materials, and enable molecular systems with applications in medicine, catalysis, and smart materials.

\subsection*{Selected Publications}

"Light-Driven Monodirectional Molecular Rotor" (Nature, 1999)

\clearpage

\section*{Chapter Summary}

The 2016 prize recognized the design and synthesis of molecular machines: Sauvage built interlocked catenanes, Stoddart constructed rotaxane shuttles, and Feringa made light-driven molecular motors. Together they showed that molecules can be engineered to perform controlled mechanical movements.
